\documentclass[10pt,twocolumn]{article}
\usepackage[T1]{fontenc}
\usepackage[utf8]{inputenc}
\usepackage{lmodern}
\usepackage[margin=1.8cm,top=2.4cm,bottom=2cm,columnsep=0.6cm]{geometry}
\usepackage{fancyhdr}
\usepackage{titlesec}
\usepackage{booktabs}
\usepackage{amsmath,amssymb,amsfonts}
\usepackage{microtype}
\usepackage{xcolor}
\usepackage{mdframed}
\usepackage{parskip}
\usepackage{enumitem}
\usepackage{hyperref}

\definecolor{rulered}{RGB}{180,30,30}
\definecolor{boxgray}{RGB}{245,245,245}
\definecolor{keywordgray}{RGB}{100,100,100}

\pagestyle{fancy}
\fancyhf{}
\fancyhead[L]{\textsc{\small Claude Mag}}
\fancyhead[C]{\small\textit{Machine Learning Research Digest}}
\fancyhead[R]{\small 2026-10-26}
\fancyfoot[C]{\small\thepage}
\renewcommand{\headrulewidth}{0.8pt}

\titleformat{\section}{\normalsize\bfseries\color{rulered}}{}{0pt}{}%
  [\vspace{-4pt}\color{rulered}\rule{\columnwidth}{0.4pt}]
\titlespacing{\section}{0pt}{8pt}{4pt}

\hypersetup{colorlinks=true,urlcolor=rulered,linkcolor=black}

\begin{document}
\setlength{\parindent}{0pt}
\setlength{\parskip}{4pt}

{\small\color{keywordgray}\textbf{Keywords:}
  safe reinforcement learning \textbullet{}
  constrained MDP \textbullet{}
  representation learning \textbullet{}
  low-rank MDP}\par\vspace{4pt}

{\LARGE\bfseries\raggedright
  Primal-Dual Representation Learning for Low-Rank Constrained MDPs}\par\vspace{4pt}

{\large\itshape\raggedright
  The First Sample-Efficient Algorithms with Theoretical Guarantees for
  Continuous-State Safe RL Under Both Soft and Hard Violation Constraints}\par\vspace{6pt}

{\small\textbf{By} Chenhao Zhou, Chao Zhang, Wensong Bai, Hanbin Zhao,
  Hui Qian}\par\vspace{2pt}

{\small\color{keywordgray}NeurIPS 2026 (Poster \#151600)}
\par\vspace{2pt}\noindent{\color{rulered}\rule{\columnwidth}{0.6pt}}\vspace{6pt}

Safe reinforcement learning requires an agent to maximize return while
keeping constraint violations within a budget---a challenge that, for
continuous high-dimensional state spaces, has lacked end-to-end
theoretical guarantees when the transition model must be learned from
data. This paper introduces REP-PD-TEG and REP-PD-H-TEG, two
primal-dual algorithms that jointly learn low-rank transition
representations and satisfy safety constraints, achieving
$\tilde{O}(\sqrt{K})$ return regret (soft violation) and certified
per-episode constraint satisfaction (hard violation) respectively. The
key insight is that representation learning and dual updates can be
decoupled: the former proceeds by maximum-likelihood exploration,
independent of the safety objective, while the latter handles constraints
via an optimistic Lagrangian with a stabilized dual search.

\begin{mdframed}[backgroundcolor=boxgray,linecolor=rulered,linewidth=1pt,
    innerleftmargin=6pt,innerrightmargin=6pt,
    innertopmargin=6pt,innerbottommargin=6pt]
\textbf{\small AT A GLANCE}\par\vspace{4pt}
\begin{description}[leftmargin=0pt,labelsep=4pt,itemsep=2pt,parsep=0pt,
    font=\small\bfseries]
  \item[Problem:] Safe RL in continuous state spaces requires learning
    transition representations and satisfying constraints simultaneously;
    no prior method provided end-to-end guarantees for both.
  \item[Why it matters:] Real-world safe RL (robotics, autonomous driving)
    has continuous state spaces; tabular guarantees do not apply, and
    deep safe RL methods (CPO, PDO) lack theoretical backing.
  \item[Method:] Decoupled representation learning via ML exploration,
    followed by optimistic Lagrangian planning (REP-PD-TEG for soft
    violation) or per-episode feasibility certification (REP-PD-H-TEG).
  \item[Result:] $\tilde{O}(\sqrt{K})$ regret and $\tilde{O}(K^{3/4})$
    total samples (soft); certified per-episode zero violation (hard).
  \item[Evidence:] Continuous-state safety gridworld, constrained LQR;
    ablations on features, exploration rate, dual step size.
\end{description}
\end{mdframed}

\section{The Big Picture}

Safe RL is the problem of training policies that maximize cumulative reward
while respecting safety constraints expressed as budgets on cumulative or
per-episode costs. It is central to deploying RL in physical systems---robots
that should not fall, vehicles that should not collide, and medical dosing
agents that should not exceed toxic thresholds.

The theoretical safe RL literature has largely studied tabular settings, where
finite-sample guarantees can be derived from counting arguments over a discrete
state space. For continuous state spaces, function approximation is necessary,
and the low-rank MDP model---where the transition kernel factorizes as
$P(s' \mid s, a) = \phi(s, a)^\top \mu(s')$ for low-dimensional features
$\phi$ and $\mu$---is the natural extension that retains mathematical tractability.

Until this paper, no work had provided end-to-end sample-efficiency guarantees
for the combined problem of learning the representation $\phi$ and satisfying
a safety constraint in a continuous-state low-rank CMDP.

\section{Why Existing Approaches Fall Short}

\textbf{Tabular constrained RL (C-UCBVI, OptPess-LP):}
Provides finite-sample guarantees but requires discretizing the state space,
which is computationally intractable in high dimensions and introduces
approximation error proportional to the discretization resolution.

\textbf{Linear function approximation for CMDPs:}
Assumes the true features $\phi$ are given; cannot handle the setting where
they must be estimated from data, which introduces feature estimation error
that can violate constraint guarantees.

\textbf{Deep safe RL (CPO, PDO, RCPO):}
Achieves strong empirical performance but provides no finite-sample bounds on
regret or constraint violation; convergence guarantees require idealized
assumptions rarely met in practice.

\textbf{Unconstrained low-rank MDP representation learning:}
Techniques such as FLAMBE and REPUCB learn $\phi$ by maximum likelihood and
achieve provably near-optimal policies---but the dual variable of the safety
constraint interacts non-trivially with the UCB bonus, breaking the direct
extension.

\section{Core Mechanism}

The key architectural insight is that representation learning and safety
constraint handling can be \emph{decoupled}:

\begin{enumerate}[itemsep=2pt,parsep=0pt]
  \item \textbf{Representation learning:} Run a decaying epsilon-greedy
    exploration policy to collect transition data; fit $\hat{\phi}$ by
    maximum-likelihood. This is independent of the safety objective.
  \item \textbf{Primal-dual planning:} Use $\hat{\phi}$ to construct UCB-style
    optimistic estimates of both the return $\hat{V}^r_\pi$ and the constraint
    utility $\hat{V}^c_\pi$. Plan with the optimistic Lagrangian:
    $\max_\pi \min_\lambda \hat{V}^r_\pi - \lambda(\hat{V}^c_\pi - b)$.
    Update $\lambda$ by dual gradient ascent.
\end{enumerate}

\textbf{REP-PD-TEG} handles the soft (cumulative) violation case with
standard UCB bonuses and gradient ascent on the dual.

\textbf{REP-PD-H-TEG} handles the hard (per-episode certified) violation case
by adding a lower-confidence constraint certificate: before deploying any policy,
the algorithm verifies that the pessimistic constraint estimate $\underline{V}^c_\pi$
is below the budget $b$, preventing episodes that violate the constraint even in
the worst case.

\section{Technical Formulation}

The low-rank CMDP structure assumes:
\[
P(s' \mid s, a) = \phi(s,a)^\top \mu(s'), \quad \phi: \mathcal{S}\times\mathcal{A} \to \mathbb{R}^d
\]
The constrained optimization problem is:
\[
\max_\pi V^r_\pi \quad \text{s.t.} \quad V^c_\pi \leq b
\]
The Lagrangian relaxation is:
\[
\mathcal{L}(\pi, \lambda) = V^r_\pi - \lambda(V^c_\pi - b), \quad \lambda \geq 0
\]
Maximum-likelihood feature learning from an epsilon-greedy exploration policy
with budget $T_0$ achieves feature error $\|\hat{\phi} - \phi^*\|_\infty \leq \epsilon_\phi$
with high probability for $T_0 = \tilde{O}(d^2/\epsilon_\phi^2)$.

With the learned features, the optimistic return and constraint estimates are:
\[
\hat{V}^r_\pi = V^r_\pi(\hat{\phi}) + \underbrace{\beta_r \cdot b^\phi_\pi}_{\text{UCB bonus}},
\quad
\hat{V}^c_\pi = V^c_\pi(\hat{\phi}) + \beta_c \cdot b^\phi_\pi
\]
The main theoretical results for REP-PD-TEG are:
\[
\text{Return regret} = \tilde{O}(\sqrt{K}), \quad
\text{Total samples} = \tilde{O}(K^{3/4})
\]
For REP-PD-H-TEG, the per-episode hard-violation guarantee holds with
probability $1-\delta$ for all $K$ episodes simultaneously.

\section{Learning or Inference Procedure}

\textbf{REP-PD-TEG:}
\begin{enumerate}[itemsep=2pt,parsep=0pt]
  \item Collect $T_0$ episodes of epsilon-greedy data; fit $\hat{\phi}$ by MLE.
  \item For each episode $k$: compute optimistic $\hat{V}^r$, $\hat{V}^c$
    using $\hat{\phi}$ and UCB bonuses; solve the optimistic Lagrangian for
    $\pi_k$; deploy $\pi_k$; update $\lambda$ by gradient ascent:
    $\lambda \leftarrow \lambda + \eta_\lambda (V^c_{\pi_k} - b)$.
\end{enumerate}

\textbf{REP-PD-H-TEG:}
\begin{enumerate}[itemsep=2pt,parsep=0pt]
  \item Steps 1--2 as above, but before deploying $\pi_k$, compute the
    pessimistic lower bound $\underline{V}^c_{\pi_k}$ using a lower-confidence
    interval on $\hat{V}^c$.
  \item If $\underline{V}^c_{\pi_k} \leq b$: deploy $\pi_k$ (certified safe).
  \item Otherwise: invoke the stabilized dual search to find a feasible policy
    (shrink dual step size if the certificate is tight).
\end{enumerate}

\section{What the Guarantee Says}

\textbf{Soft violation (REP-PD-TEG):}
The cumulative return regret over $K$ episodes is $\text{Reg}_K = \tilde{O}(\sqrt{K})$,
and the total cumulative constraint violation is bounded by $\tilde{O}(\sqrt{K})$.
Total episodes including exploration: $\tilde{O}(K^{3/4})$.

\textbf{Hard violation (REP-PD-H-TEG):}
With probability $1-\delta$, the deployed policy satisfies $V^c_{\pi_k} \leq b$
for all $k = 1, \ldots, K$ simultaneously. Return regret is slightly worse than
the soft case due to the conservative certification step but remains
$\text{poly}(K)$-sublinear.

\textbf{First result of its kind:}
The authors claim these are the first finite-sample guarantees for
representation learning in low-rank CMDPs with continuous state spaces
under both soft and certified hard violation settings.

\section{Experimental Findings}

\begin{itemize}[itemsep=2pt,parsep=0pt]
  \item \textbf{Continuous safety gridworld:} REP-PD-TEG keeps cumulative
    constraint violations below budget; REP-PD-H-TEG achieves zero per-episode
    violations after an initial feature-learning phase.
  \item \textbf{Constrained LQR:} Both algorithms outperform CPO and PDO
    (which use fixed features) when the provided features are incorrect;
    performance matches CPO with correct features.
  \item \textbf{Benefit of learned features:} Methods with fixed (incorrect)
    features consistently violate the constraint budget; the adaptive feature
    learning of REP-PD-TEG removes this failure mode.
\end{itemize}

\section{Ablations and Interpretation}

\textbf{Fixed vs.\ learned features:}
Ground-truth features in the primal-dual step achieve near-optimal performance;
incorrect fixed features fail to satisfy constraints and achieve poor return,
confirming the necessity of end-to-end representation learning.

\textbf{Epsilon-greedy schedule:}
A decaying schedule is critical; a constant rate leads to over-exploration and
slower convergence; an overly fast decay leads to poor feature quality and
constraint violations.

\textbf{Dual step size stability (REP-PD-H-TEG):}
The stabilized dual search (reduced step size near the constraint boundary)
is essential for hard-violation safety; naive gradient ascent on the dual
oscillates and produces intermittent per-episode violations.

\section*{Reference}

Chenhao Zhou, Chao Zhang, Wensong Bai, Hanbin Zhao, Hui Qian.
\textbf{Primal-Dual Representation Learning for Low-Rank Constrained MDPs}.
NeurIPS 2026 (Poster \#151600).
\url{https://neurips.cc/virtual/2026/poster/151600}

\end{document}
